\chapter{Wounds}
\index{Wounds}
\section*{Definition and Overview}
A wound is a break or damage to the skin or deeper tissue caused by injury, surgery, pressure, vascular disease, infection, or inflammation. Wounds may be acute or chronic and may involve nerves, tendons, vessels, bone, or internal organs.
\section*{Epidemiology}
Common wounds include cuts, grazes, punctures, bites, burns, surgical wounds, pressure injuries, and diabetic or vascular ulcers. Healing is slower with infection, poor blood supply, diabetes, smoking, malnutrition, immune suppression, oedema, and ongoing pressure.
\section*{Aetiology and Pathophysiology}
Healing involves haemostasis, inflammation, tissue formation, and remodelling. Contamination, dead tissue, foreign material, poor perfusion, repeated trauma, and excessive tension can interrupt this process and lead to infection or chronic non-healing.
\section*{Clinical Features}
Assess wound location, size, depth, bleeding, contamination, pain, tissue viability, sensation, movement, pulses, and function. Redness spreading from the wound, increasing pain, pus, fever, or wound separation suggests complication.
\section*{Differential Diagnosis}
Consider laceration, puncture, crush injury, bite, burn, pressure injury, venous or arterial ulcer, diabetic foot ulcer, skin cancer, pyoderma gangrenosum, vasculitis, and infected eczema.
\section*{When to Seek Medical Care}
Urgent assessment is needed for uncontrolled bleeding, deep or gaping wounds, injuries involving the face, eye, hand, genitalia, joint, tendon, nerve, or bone, retained foreign body, animal or human bites, puncture wounds, severe burns, loss of sensation or movement, or a wound in a person with diabetes or poor circulation. Emergency care is required for shock, severe infection, or rapidly spreading redness with systemic illness.
\section*{Diagnosis and Investigations}
Most minor wounds need examination only. X-ray or other imaging may be needed for foreign bodies, fracture, deep infection, or penetrating injury. Wound cultures are reserved for clinically infected or unusual wounds. Assess tetanus immunisation and, after relevant exposure, blood-borne virus risk.
\section*{Management}
Apply direct pressure to bleeding, rinse minor wounds with clean running water, remove only easily accessible debris, and cover with a clean dressing. Do not routinely use hydrogen peroxide, alcohol, or harsh antiseptics inside wounds. Clinicians may close suitable wounds, debride dead tissue, prescribe antibiotics for selected infections or high-risk injuries, and update tetanus protection. Chronic wounds require pressure relief, vascular assessment, infection management, nutrition, and treatment of the underlying disease.
\section*{Complications}
Complications include infection, abscess, cellulitis, tetanus, scarring, delayed healing, nerve or tendon damage, tissue loss, osteomyelitis, sepsis, and chronic ulceration.
\section*{Prognosis and Outcomes}
Simple clean wounds often heal within days to weeks. Healing time depends on depth, location, blood supply, infection, underlying disease, and treatment. Persistent wounds require reassessment rather than repeated self-treatment.
\section*{Prevention and Self-Care}
Use protective equipment, control diabetes and vascular disease, stop smoking, maintain nutrition, prevent pressure and falls, keep dressings clean and dry as advised, and follow up when instructed.
\section*{Sources and Further Reading}
\begin{itemize}
    \item NHS cuts and grazes: \url{https://www.nhs.uk/conditions/cuts-and-grazes/}
    \item World Health Organization wound care: \url{https://www.who.int/teams/integrated-health-services/clinical-services-and-systems/surgical-care}
\end{itemize}
